% ---------- Оформление по ГОСТ ----------
% Поля: слева 30, справа 15, сверху/снизу 20 мм
\usepackage[a4paper, left=30mm, right=15mm, top=20mm, bottom=20mm]{geometry}

\setcounter{section}{-1}

% Интервал 1,5; абзацный отступ 1,25 см, в том числе после заголовков
\usepackage{setspace}
\onehalfspacing
\usepackage{indentfirst}
\setlength{\parindent}{1.25cm}

% Заголовки: разделы --- ЗАГЛАВНЫМИ, жирным; подразделы --- жирным
\usepackage{titlesec}
\titleformat{\section}{\normalfont\bfseries}{\thesection}{1em}{\MakeUppercase}
\titleformat{\subsection}{\normalfont\bfseries}{\thesubsection}{1em}{}

% Подписи: «Рисунок 1 – Название», «Таблица 1 – Название»
% \usepackage{caption}
% \captionsetup{labelsep=endash, font=normalsize}
% \captionsetup[table]{justification=raggedright, singlelinecheck=off}
% \captionsetup[figure]{justification=centering}
% \addto\captionsrussian{%
%   \renewcommand{\figurename}{Рисунок}%
%   \renewcommand{\tablename}{Таблица}}
% \renewcommand{\lstlistingname}{Листинг}

\usepackage{amsmath, amsthm}
\theoremstyle{definition}
\newtheorem{definition}{Опр.}[section]
\newtheorem{example}{Пример}[section]

\theoremstyle{plain}
\newtheorem{theorem}{Теорема}[section]
\newtheorem{lemma}{Лемма}[section]
\newtheorem{statement}{Утв.}[section]
\newtheorem{note}{Замечание}[section]

\usepackage[most]{tcolorbox}
\newtcolorbox{commentary}[1][]{
    enhanced,
    colback=gray!5,                 % Серый ультра-легкий фон
    colframe=gray!40,               % Тонкая рамка чуть темнее фона
    boxrule=0.5pt,                  % Толщина рамки
    sharp corners,                  % Квадратные углы (строго по ГОСТ)
    before skip=10pt,               % Отступ сверху от текста
    after skip=10pt,                % Отступ снизу до текста
    fontupper=\small\fontspec{PT Sans}, % Мелкий, тонкий рубленый шрифт без засечек
    left=10pt, right=10pt,          % Внутренние отступы слева/справа
    top=8pt, bottom=8pt,            % Внутренние отступы сверху/снизу
    #1
}